
\def\PUDcodigo{ECA227}

\def\PUDcodacad{04507.45}

\def\PUDdisciplina{Sistemas Digitais de Controle Distribuído}

\def\PUDsemestre{S9}

\def\PUDhoras{80}

%NOVOS ---------------------------------------------------
\def\PUDhorasTeoria{60}
\def\PUDhorasPratica{20}

\def\PUDinterECA{ 
\hyperref[PUD:ECA209]{Linguagem de Programação (04507.14)}

\hyperref[PUD:ECA219]{Instrumentação (04507.25)}

\hyperref[PUD:ECA225]{Dispositivos Periféricos (04507.42)} 
}

\def\PUDatuaisECA{
- Noções de Internet das coisas (IoT) e aplicações
}

\def\PUDrecursos{Projetor multimídia; Quadro branco e pincel; Aulas em laboratório.}

%----------------------------------------------------------

\def\PUDcreditos{4}

\def\PUDprerequisito{---}

\def\PUDpreacad{---}

\def\PUDobjetivos{Compreender os conceitos de sistema distribuído;  Implementar sistemas de controle supervisório baseados em redes de dispositivos de controle e sistemas SCADA para sistemas de manufatura e controle de processos.  
Conhecer os conceitos de redes de dados, estruturas e camadas de comunicação; 
Conhecer estruturas, protocolos e gerenciamento de redes industriais utilizadas em sistemas de supervisão; 
Planejar e implementar um sistema de comunicação de dados em rede de dispositivos.}

\def\PUDementa{Introdução ao Controle distribuído; 
Introdução ao SCADA; 
Desenvolvimento de Aplicativos utilizando controladores em rede;  Sistemas SCADA; 
Desenvolvimento de Aplicativos SCADA; 
Conceitos de redes de computadores;  Modelo OSI/ISO;  Topologias de rede;  Redes industriais de controle e supervisão. 
Conceitos  de redes de computadores;  topologias;  modelo OSI/ISO;   transmissão serial;   protocolos industriais;  Telemetria convencional a 2 fios/4 fios;  Redes industriais;  Estrutura de redes industriais: Fieldbus, Devicebus e sensorbus;  Protocolos de comunicação de redes industriais; Gerenciamento de redes industriais;  Manutenção de redes industriais; Introdução aos sistemas supervisórios.}

\def\PUDconteudo{ & Introdução ao Sistema de Aquisição de Dados e Controle Supervisório; 
\\ 
 & Características dos sistemas SCADA; 
\\ 
 & Interface homem-máquina gráfica; 
\\ 
 & Desenvolvimento de aplicativos SCADA e supervisórios; 
\\ 
 & Conceito de redes: Redes de computadores;  Redes industriais; Redes industriais;  \\ 
 & Telemetria convencional a 2 fios/4 fios:  Barramentos;  Comunicação Simplex, Half-duplex e Duplex; \\ 
 & Modelo OSI/ISO: Camada Física;  Camada Enlace de dados; Camada de Rede; Camada de Transporte; Camada de Sessão; Camada de Apresentação; Camada de Aplicação; \\ 
 & Modelos de Redes Industriais: Introdução ao Processo Industrial; Descrição do Processo; Barramento de Campo;  Cabeamento;  Infra-estrutura; \\ 
 
 & Estrutura de Redes Industriais: Fieldbus,Devicebus e sensorbus;  Redes de Sensores; Redes de Controle; Noções de Internet das coisas (IoT) e aplicações.
\\ 
 & Protocolos de rede: Protocolos Abertos; Protocolos Fechados; Profibus, Modbus, Fieldbus, Hart. }

\def\PUDmetodo{Aulas expositórias que podem ser teóricas e/ou práticas, onde as práticas no laboratório serão marcadas no decorrer da disciplina.}

\def\PUDavalia{A avaliação será feita com aplicação de provas teóricas e/ou práticas, além da possibilidade de inclusão de trabalhos e seminários no decorrer da disciplina.}

\def\PUDbibbasica{ & \href{www.biblioteca.ifce.edu.br/index.asp?codigo_sophia=24170}{Albuquerque}, Pedro Urbano Braga;  Alexandria, Auzuir Ripardo.  Redes industriais: aplicações em sistemas digitais de controle distribuído: protocolos industriais,  aplicações SCADA.  2º ed.  ISBN: 9788599823118.
\\ 
 & \href{www.biblioteca.ifce.edu.br/index.asp?codigo_sophia=24230}{Campos}, Mário Cesar M.  Massa de;  Teixeira, Herbert C.  G.  Controles típicos de equipamentos e processos industriais.  São Paulo, SP: Edgard Blücher, 2006.  396p.  ISBN: 9788521203988. 
\\ 
 &   \href{biblioteca.ifce.edu.br/index.asp?codigo_sophia=40142}{Young}, Paul H. Técnicas de comunicação eletrônica. 5. ed. São Paulo: Pearson Prentice Hall, 2014. 687 p. ISBN 9788576050490. }

\def\PUDbibcomplementar{ &  \href{www.biblioteca.ifce.edu.br/index.asp?codigo_sophia=24297}{Alves}, J. L. L.  Instrumentação, Controle e Automação de Processos. Rio de Janeiro, RJ: LTC, 2005.  270p.  ISBN: 9788521614425.
\\ 
 &  \href{www.biblioteca.ifce.edu.br/index.asp?codigo_sophia=24327}{Natale}, Ferdinando.  Automação industrial.  10º ed. São Paulo, SP: Érica, 2011.  252p.  (Série Brasileira de Tecnologia).  ISBN: 9788571947078.
\\ 
 &  \href{www.biblioteca.ifce.edu.br/index.asp?codigo_sophia=65382}{Capelli}, Alexandre.  Automação industrial: controle do movimento e processos contínuos.  3º ed.  São Paulo, SP: Érica, 2013.  236p.  ISBN: 9788536501178.
\\
 &  \href{www.biblioteca.ifce.edu.br/index.asp?codigo_sophia=52512}{Brito}, Fábio Timbó.  Protocolos de comunicação.  Curitiba, PR: Editora do Livro Técnico, 2015.  208p.  ISBN: 9788563687807.
\\
 &  \href{www.biblioteca.ifce.edu.br/index.asp?codigo_sophia=40178}{Pinheiro}, José Maurício.  Infra-estrutura elétrica para rede de computadores.  Rio de Janeiro, RJ: Ciência Moderna, 2008.  281p.  ISBN: 9788573936865. }
%\\ 
 %&  MORAES, C.  C. , Engenharia de Automação Industrial, 2a Edição, Rio de Janeiro, Editora LTC, 2007.  ISBN 9788521615323
%\\ 
 %& TANENBAUM, A.  S. , Redes de Computadores, 1a Edição, Editora Campus.  2003.  ISBN 9788535211856
%\\ 
 %&  SOARES, L.  F. , Redes de Computadores, 2a Edição, Editora Campus.  1995.  728p.  ISBN 9788570019981
%\\ 
 %&   RAPPAPORT, T.  S.  Comunicações sem Fio: Princípios e práticas.  2a.  Edição.  São Paulo, Pearson Prentice Hall, 2009.  409 p.  ISBN 9788576051985. }

\def\PUDautor{Geraldo Luis Bezerra Ramalho}

\def\PUDdata{2013-05-20}

\def\PUDrevisao{2}

\def\PUDrevisor{Fábio Timbó Brito}

\def\PUDdatarevisao{2019-05-15}

\PUDcorpo
